\chapter{Conclusion and Future Work}\label{ch:conclusion}


This dissertation set out to show that useful quantum computation in the near term is a property of a co-designed stack rather than of a quantum processor alone, and to make concrete progress at each layer of that stack. This chapter returns to the five research questions posed in \cref{ch:intro}, states what the preceding chapters have established and with what evidence, draws out the design principles that recur across layers, records the limitations that remain, and lays out the research program that follows.

\section{Answers to the Research Questions}\label{sec:conclusion:answers}

\subsection{RQ1: Scale}\label{sec:conclusion:rq1}

\Cref{ch:qgear} asked how far classical HPC can push exact circuit simulation and how that capability can be made usable without rewriting programs. The answer given by \qgear{} has two parts. Technically, exact state-vector simulation can be pushed to 42 qubits on 1,024 A100 GPUs of Perlmutter with close to full GPU utilization, and the path from a laptop-scale \qiskit{} circuit to that scale is a constant-time gate-by-gate transformation into a \cudaq{} kernel plus a containerized, Slurm-native deployment~\cite{guo2025qgear}. Practically, the speed-ups that matter to a working researcher are the ones that change what experiments are feasible in a day: a 34-qubit random circuit that occupies a CPU node for a day completes in about a minute on four GPUs, a 20-qubit quantum Fourier transform runs one to two orders of magnitude faster than in \qiskit{} Aer, and an 18-qubit variational eigensolver loop drops from ten minutes to four seconds. The framework was subsequently used as the validation and reward-evaluation engine of \cref{ch:nisq,ch:qml,ch:synthesis}, which is the strongest evidence that the usability goal was met. The limit is also clear: the memory wall of exact simulation stands at roughly 42--44 qubits even for a leadership-class allocation, so beyond that scale simulation must become approximate (tensor networks, Clifford-plus-magic decompositions) or give way to the hardware.

\subsection{RQ2: Data}\label{sec:conclusion:rq2}

\Cref{ch:qml} asked how classical data should enter and leave a quantum processor so that learning models are shot-efficient, hardware-compatible and trainable. The four primitives developed there share one answer: load many values at once into an address--data register with a uniformly controlled rotation, read many expectation values at once from the same circuit, keep the nonlinearity on the classical side, and let a hardware-aware compiler decide how the circuit meets the device. \qpie{} demonstrated that parallel exchange between pretrained classical parameters and variational circuits improves accuracy over classical transfer learning on five benchmark tasks and speeds convergence on stochastic time series~\cite{guo2025qpie}. \vqt{} demonstrated that the attention matrix of a transformer can be estimated by a vectorized quantum dot product with logarithmic circuit complexity in the batch--sequence product, with an RMSE of 0.059 on \texttt{ibm\_kingston}, and that the resulting language model is competitive with prior quantum sequence models on the Brown corpus~\cite{guo2025vqt}. \shardq{} demonstrated that a deterministic, distance-ordered cut selection with global reconstruction lowers encoding error by about 15\% on Heron devices and lifts reconstruction fidelity from 0.79 to 0.90 for seven-address-qubit images~\cite{guo2025shardq}. \nnqa{} demonstrated that a trained polynomial network can be compiled into exact quantum arithmetic that reaches above 99.5\% accuracy for polynomials up to degree 35 on both superconducting and trapped-ion hardware, with a resource footprint an order of magnitude below variational alternatives~\cite{guo2026nnqa}.

\subsection{RQ3: Noise}\label{sec:conclusion:rq3}

\Cref{ch:nisq} asked how optimization algorithms on NISQ devices can exploit problem and hardware structure instead of black-box parameter search. \deal{} showed that a QUBO instance carries enough structure to initialize QAOA angles directly through qubit-prioritization normalization, that calibration data can drive a dynamic physical mapping, and that a lightweight digital zero-noise extrapolation can be folded into the same pipeline; on IBM Heron processors this raised the ground-state success rate of 20-qubit, 10-layer circuits by 14.81\% (\texttt{ibm\_torino}) and 7.40\% (\texttt{ibm\_marrakesh}) over vanilla QAOA, and let seven layers do the work of nine~\cite{guo2025deal}. \qawa{} approached the same question from the data side: a classical-data-traceable oracle whose two-qubit depth is linear in the qubit count, together with classical post-processing of mid-circuit measurements, yields solutions whose quality can be verified in polynomial time, and on a 156-qubit device produced a portfolio 5.4\% better than the baseline and 8.1\% above the noise floor~\cite{guo2025qawa}. The general answer is that structure should be moved out of the optimizer and into the ansatz, the initialization and the compilation, where it costs no shots.

\subsection{RQ4: Synthesis}\label{sec:conclusion:rq4}

\Cref{ch:synthesis} asked whether learned, HPC-scale synthesis can produce circuits that respect fault-tolerant resource budgets, and what honest cost accounting says about end-to-end quantum advantage. \rubriq{} answered the first half: a code-generating language model fine-tuned with group relative policy optimization against a programmatic rubric, with GPU simulation inside the reward loop, achieves a mean T-gate compression of 3.31$\times$ against 2.05$\times$ for sparse-reward reinforcement learning, keeps hardware-constraint violations below 1\%, and converges two to three times faster; the dense, decomposable reward is what makes the difference~\cite{guo2026rubriq}. The nonlinear-wave study answered the second half with a negative result that is stated quantitatively: in a split-step solver that must measure the field to evaluate the nonlinearity, the per-step quantum cost, measured as depth times shots, grows faster than the classical cost with grid size, so an end-to-end advantage requires a coherent nonlinear update that does not read out the field~\cite{guo2026nonlinearwaves}. Both halves rely on the same HPC substrate as \cref{ch:qgear}, and together they set the bar that fault-tolerant compilation must clear.

\subsection{RQ5: Trust}\label{sec:conclusion:rq5}

\Cref{ch:qec,ch:pqc} asked how error-corrected computation on shared cloud QPUs can be made auditable and secure, and how that connects to post-quantum cryptography standards. The audit of \cref{ch:qec} introduced a bounded cloud fault model, the notion of accepted logical disturbance and its exact expectation under Haar-random encoders, and a polynomial-cost seeded Clifford ensemble that reseeds the encoder per run; for eight-qubit encoders the reseeded ensemble reduces mean accepted disturbance from 0.150 to 0.020, an 86.7\% reduction attributable to syndrome-based rejection, while the fixed \code{5}{1}{3} code corrects every tested weight-one fault and gate-level Stim simulations up to seventeen qubits confirm the scaling~\cite{guo2026qecaudit}. \Cref{ch:pqc} then showed that the seed for that encoder, the keys that protect job submission, and the signatures that authenticate results should all come from a post-quantum cryptographic module designed to the FIPS 140-3 requirements, and described the Light Rider QPC SDK as such a module with a discovery service for approved algorithm suites~\cite{guo2026lightriderqpc}. The answer to RQ5 is therefore that trust in the quantum stack is a chain: validated randomness feeds reseeded encoders, reseeded encoders bound what a co-tenant adversary can do, and post-quantum key establishment and signatures protect everything that moves between the classical and quantum sides.

\section{Cross-Cutting Principles}\label{sec:conclusion:principles}

Four principles recur across the chapters and are, in the candidate's view, the durable contribution of the dissertation beyond its individual results.

The first is \emph{move work off the QPU and out of the shot budget}. \deal{} moves parameter search into initialization and compilation; \qawa{} and \vqt{} move nonlinear post-processing to the classical side; \nnqa{} moves training entirely off the device and uses it only for exact arithmetic; \rubriq{} moves the search for good circuits into a classical policy trained against a simulator. In every case the quantum processor is asked to do the one thing it does that the classical side cannot, and nothing else.

The second is \emph{vectorize the interface}. Address--data encodings and parallel readout appear in \cref{ch:qgear} as the \qcrank{} benchmark, in \cref{ch:qml} as the foundation of all four primitives, in \cref{ch:nisq} as the data-traceable oracle, and in \cref{ch:synthesis} as the reload step of the split-step solver. Whatever the application, the cost of the classical--quantum interface is set by how many values cross it per circuit execution, and vectorization is the lever.

The third is \emph{let calibration data drive compilation}. \deal's dynamic mapping and \rubriq's hardware-compliance term consume device calibration and topology as inputs to compilation rather than as an afterthought, and \shardq{} transpiles its sub-circuits against the calibrated coupling map of the target. The devices used in this dissertation changed their error maps daily, and methods that read the map outperformed methods that did not.

The fourth is \emph{account honestly and audit adversarially}. The nonlinear-wave study reports where the quantum cost overtakes the classical one; the QEC audit assumes an adversary rather than benign noise; and \cref{ch:pqc} distinguishes what has been validated from what has been designed. A stack that is to be trusted with computations of value must be assessed against the costs it will actually incur and the adversaries it will actually face.

\section{Limitations}\label{sec:conclusion:limitations}

The limitations recorded in each chapter fall into three groups. The first is scale on the quantum side. Hardware experiments used at most 36 qubits and a few hundred two-qubit gates; the algorithms of \cref{ch:nisq,ch:qml} were shown to help at that scale, and simulation extends the evidence to 42 qubits, but whether the observed gains persist at hundreds of logical qubits is an open empirical question. The second is model scope. The QEC audit uses weight-one Pauli faults and small codes; the cost model of the nonlinear-wave study assumes a particular loader and a particular polynomial nonlinearity; and \rubriq's rubric encodes one particular view of fault-tolerant cost centred on T-count. Each model was chosen to be tractable and defensible, and each could be broadened. The third is validation status. Several results are preprints under review at the time of writing, and the cryptographic module of \cref{ch:pqc} is designed to a standard rather than certified against it. These facts are stated where they arise, and the artifact appendix (\cref{app:repro}) records the status of each work.

\section{Future Work}\label{sec:conclusion:future}

The research program that follows from this dissertation has one item per layer and one that spans them.

At the \emph{scale} layer, the natural extension of \qgear{} is approximate simulation beyond the memory wall: matrix-product-state and tree-tensor-network backends whose bond dimension is chosen automatically from the circuit's entanglement structure, and Clifford-plus-magic decompositions that exploit the low T-count of the circuits \rubriq{} produces. A second extension is hardware-in-the-loop simulation in which the noise model is refreshed from the device's live calibration, so that the simulator predicts hardware results rather than ideal ones.

At the \emph{data} layer, the four primitives of \cref{ch:qml} should be unified into a single vectorized-encoding library with a common compiler pass built on \shardq{}, and \vqt{} should be scaled to sequence lengths and vocabularies where the logarithmic circuit complexity of the vectorized dot product becomes the dominant advantage. \nnqa's exact arithmetic suggests a broader program of compiling classically trained models into quantum arithmetic blocks, including the coherent nonlinear update that \cref{ch:synthesis} identifies as the missing piece for quantum simulation of nonlinear dynamics.

At the \emph{algorithms} layer, \deal{} and \qawa{} should be combined: the direct initialization and dynamic mapping of \deal{} with the data-traceable oracle and verifiable post-processing of \qawa{}, applied to constrained problems with penalty structure. Their integration with error-corrected execution, in which the logical-qubit error map replaces the physical one, is the longer-term goal.

At the \emph{synthesis} layer, \rubriq's rubric should be extended from T-count to full fault-tolerant resource estimates (logical qubits, code distance, magic-state factories) so that the policy optimizes the quantity a fault-tolerant machine will actually charge for, and the simulator in the reward loop should be replaced, for larger circuits, by the approximate backends proposed above.

At the \emph{trust} layer, the QEC audit should be extended to correlated and higher-weight faults, to surface and bivariate-bicycle codes, and to an adversary who observes syndromes; and the QPC SDK should proceed through formal FIPS 140-3 validation with the hybrid key-establishment and signed-result protocols of \cref{ch:pqc} exercised against live cloud QPU services.

Spanning all layers, the most important next step is an end-to-end demonstration: a single workload developed in simulation with \qgear, compiled with \shardq{} and \rubriq, executed with \deal-style initialization on a cloud QPU under a reseeded error-correcting encoder whose seeds and job channels are protected by the QPC SDK, with the full cost accounted for. That demonstration would test the thesis of this dissertation directly, and the components to build it now exist.

\section{Concluding Remarks}\label{sec:conclusion:remarks}

The quantum processor is the most photographed part of a quantum computer, but it is not where most of the computation happens, and it will not be for some time. This dissertation has argued, and shown across simulation, algorithms, data encoding, synthesis and security, that treating the classical HPC around the QPU as a first-class object of research produces gains at every layer and that those gains compound. The stack is the machine.
